% ============================================================
% Weekly Meeting Report — Weeks 15--16
% Topic: Environment Setup & Dataset Audit
% ============================================================

\section*{Weekly Progress Report — Weeks 15--16}
\addcontentsline{toc}{section}{Weekly Progress Report — Weeks 15--16}

\begin{tabular}{ll}
    Date      & \today \\
    Student   & Khac Duc Giang Nguyen \\
    Supervisor & Dr.\ Seyed Sahand Mohammadi Ziabari \\
    Phase     & Environment Setup \& Dataset Audit (Weeks 15--16) \\
\end{tabular}

\vspace{1em}

% ------------------------------------------------------------
\subsection*{1. Overview}

This report documents the work completed during Weeks 15--16 of the
thesis project. The primary objectives were to establish a reproducible
working environment on the UvA Snellius supercomputer and to conduct a
thorough audit of all four datasets prescribed in the thesis proposal.
No model training has begun; the focus of this phase was entirely on
infrastructure and data understanding, in line with the Environment Setup milestone.

% ------------------------------------------------------------
\subsection*{2. Snellius Environment Setup}

A dedicated Conda environment (\texttt{uav\_master}, Python 3.9) was
created on the Snellius HPC cluster under the project directory
\texttt{/projects/prjs2041/}. Due to the constraints of the
supercomputer architecture, a hybrid installation strategy was adopted:
Conda was used only to create the Python sandbox, while \texttt{pip}
(with \texttt{--default-timeout=1000}) was used to install all packages,
bypassing the Conda dependency solver which is incompatible with login
node resource limits.

\subsubsection*{2.1 Installed Libraries}

\begin{itemize}
    \item Core deep learning (CUDA 11.8): \texttt{torch}
          (2.7.1+cu118), \texttt{torchvision} (0.22.1+cu118),
          \texttt{torchaudio} (2.7.1+cu118).
    \item Computer vision and data processing:
          \texttt{opencv-python}, \texttt{matplotlib},
          \texttt{numpy} (2.0.2), \texttt{scipy}, \texttt{pandas}.
    \item Utilities: \texttt{wget}, \texttt{shapely},
          \texttt{PyYAML}, \texttt{tqdm}, \texttt{ipykernel}.
\end{itemize}

\subsubsection*{2.2 Roadblocks Encountered}

Setting up the environment required resolving four non-trivial
infrastructure problems specific to the Snellius HPC architecture.

Problem 1 — \texttt{conda install} froze on the login node.
Conda's dependency solver requires significant RAM and CPU. Login nodes
(\texttt{int5}/\texttt{int6}) throttle both heavily, causing the process
to freeze without error. Fixed by switching to \texttt{pip install} with
an extended timeout.

Problem 2 — Lmod module loading failure.
Snellius uses a hierarchical module system where Conda is hidden until
the correct software stack year is loaded. Fixed by running
\texttt{module load 2023} before \texttt{module load Miniconda3/23.5.2-0}.

Problem 3 — VS Code module hijacking (\texttt{ModuleNotFoundError: torch}).
Even with the \texttt{uav\_master} environment active in the terminal,
the system's own aliases intercepted the \texttt{python} command and
pointed it to a different interpreter. Fixed by always calling the
absolute path of the Conda Python binary:
\begin{center}
\texttt{/home/knguyen1/.conda/envs/uav\_master/bin/python}
\end{center}

Problem 4 — VS Code UI glitches and Pylance \texttt{SIGTERM} crashes.
The Open OnDemand VS Code interface failed to auto-detect environments
stored in hidden Linux directories (\texttt{.conda}). Fixed by creating
\texttt{.vscode/settings.json} in the project folder, explicitly setting
\texttt{python.defaultInterpreterPath} to the absolute Conda binary path.

\subsubsection*{2.3 Standard Session Startup Procedure}

Every new VS Code session on Snellius requires the following sequence
before running any code:

\begin{enumerate}
    \item Launch VS Code via the SURF Open OnDemand portal
          (\texttt{rome}/\texttt{genoa} partition, 4 cores, 16\,GB RAM
          for data work; \texttt{gpu\_a100} for training).
    \item Open workspace \texttt{/projects/prjs2041/uav\_code/} and a
          new terminal.
    \item Run the startup sequence:
\begin{verbatim}
module load 2023
module load Miniconda3/23.5.2-0
source activate uav_master
\end{verbatim}
    \item Run all scripts via the absolute path rather than the
          \texttt{python} alias:
\begin{verbatim}
/home/knguyen1/.conda/envs/uav_master/bin/python script.py
\end{verbatim}
\end{enumerate}

% ------------------------------------------------------------
\subsection*{3. Dataset Extraction}

All four raw dataset archives were transferred and extracted to
\texttt{/projects/prjs2041/datasets/}. The CST Anti-UAV dataset required
special handling: the primary concatenated archive
(\texttt{combined\_archive.zip}) was found to be corrupt. Extraction was
performed instead from the original eleven-part split archive
(\texttt{CST-AntiUAV.zip} + \texttt{.z01}--\texttt{.z10}) using
\texttt{p7zip}, loaded via the Snellius module system and executed as a
dedicated SLURM batch job on the \texttt{rome} partition, yielding
246,131 files across 226 folders. SLURM job submission and monitoring
workflows are now established for future long-running training jobs.

% ------------------------------------------------------------
\subsection*{4. Dataset Storage Formats}

Table~\ref{tab:storage} summarises the storage format and modality of
each dataset. This distinction is critical for pipeline planning: two
datasets store footage as raw video files and require a frame extraction
step before any training data loader can read them, while the other two
are already stored as pre-extracted image frames and are ready to use
directly.

\begin{table}[h]
\centering
\caption{Storage format and modality per dataset.}
\label{tab:storage}
\begin{tabular}{llll}
\hline
Dataset & Storage Format & Modality & Ready for training? \\
\hline
Anti-UAV-RGBT  & MP4 video (paired RGB + IR) & RGB + Thermal & No — extraction needed \\
Anti-UAV410    & Pre-extracted frames (JPG)  & Thermal only  & Yes \\
ARD100         & MP4 video                   & RGB only      & No — extraction needed \\
CST Anti-UAV   & Pre-extracted frames (JPG)  & Thermal only  & Yes \\
\hline
\end{tabular}
\end{table}

\noindent Consequently, before any training pipeline can be assembled,
frame extraction scripts must be written for Anti-UAV-RGBT (both the
infrared and visible streams) and ARD100. This is the first concrete
engineering task scheduled for the next phase.

% ------------------------------------------------------------
\subsection*{5. Dataset Audit}

A Python audit script was written and executed across all four datasets
to verify integrity, confirm annotation formats, and compute bounding
box size distributions. The size categories follow the diagonal-length
classification defined in the thesis proposal: tiny
($d \leq 10$\,px), small ($d \in (10, 30]$\,px),
normal ($d \in (30, 50]$\,px), and large ($d > 50$\,px).

\subsubsection*{5.1 Anti-UAV-RGBT (T\textsubscript{1} Source Domain)}

\begin{itemize}
    \item Sequences: 91 \quad Frames: 85,374
    \item Annotation format: Per-sequence JSON with \texttt{exist}
          flags and \texttt{gt\_rect} bounding boxes in $[x, y, w, h]$
          format; separate files for infrared (\texttt{infrared.json}) and
          visible (\texttt{visible.json}) streams.
    \item Temporal alignment: All 91 sequences confirmed to have
          identical frame counts across IR and RGB streams — a prerequisite
          for the Teacher-Student domain adaptation framework.
    \item IR stream (Student input): 98.6\% of frames annotated
          as visible; size distribution dominated by normal
          (28.2\%) and large (68.4\%) targets.
    \item RGB stream (Teacher input): 93.4\% visible;
          98.0\% of annotated frames fall in the large category,
          reflecting the higher spatial resolution of the visible-light
          sensor relative to the IR sensor.
\end{itemize}

\subsubsection*{5.2 Anti-UAV410 (T\textsubscript{1} Target Domain)}

\begin{itemize}
    \item Sequences: 200 train / 90 val / 120 test = 410 total
    \item Annotation format: Per-sequence JSON (same schema as
          Anti-UAV-RGBT) plus global COCO-format JSON files under
          \texttt{annotations/}.
    \item Visibility: $\approx$97--98\% across all splits.
    \item Size distribution (train): tiny 2.5\%, small 27.5\%,
          normal 26.8\%, large 43.2\%. Notably more diverse in scale
          than the RGBT source domain, with a substantial small/normal
          proportion introducing an initial scale challenge even within
          the T\textsubscript{1} modality-shift task.
\end{itemize}

\subsubsection*{5.3 ARD100 (Ego-Motion Stress Test)}

\begin{itemize}
    \item Sequences: 100 \quad Frames: 202,411
    \item Annotation format: Pascal VOC XML, one file per frame,
          with bounding boxes in $[\mathit{xmin}, \mathit{ymin},
          \mathit{xmax}, \mathit{ymax}]$ format.
    \item Visibility: 98.5\%.
    \item Size distribution: Dominated by small targets
          (75.3\%), consistent with the air-to-air capture scenario where
          the target drone occupies a modest fraction of the frame.
          This confirms ARD100 as a suitable stress test for the
          Global Motion Compensation (GMC) module.
\end{itemize}

\subsubsection*{5.4 CST Anti-UAV (T\textsubscript{2} --- Extreme Scale Shift)}

\begin{itemize}
    \item Sequences: 120 train / 40 val / 60 test = 220 total
    \item Annotation format: Per-sequence \texttt{gt.txt}
          (comma-separated floating-point $[x, y, w, h]$) combined with
          \texttt{IR\_label.json} for existence flags.
    \item Visibility: 82--86\% — significantly lower than all
          other datasets, consistent with the high prevalence of Thermal
          Crossover events documented in the thesis proposal.
    \item Size distribution (train): tiny 20.5\%, small 73.6\%,
          normal 5.9\%, large $<$0.01\%. The near-complete absence of
          large targets and the 14.1\% occlusion rate constitute a
          qualitatively different detection regime from T\textsubscript{1}.
\end{itemize}

% ------------------------------------------------------------
\subsection*{6. Key Findings Relevant to Hypothesis H1}

The audit quantitatively confirms the scale-distribution shift
hypothesised in H1. Table~\ref{tab:scale_shift} summarises the
contrast between the T\textsubscript{1} target domain and the
T\textsubscript{2} domain across the two most informative splits.

\begin{table}[h]
\centering
\caption{Size distribution shift from T\textsubscript{1} to
         T\textsubscript{2} (percentage of annotated frames).}
\label{tab:scale_shift}
\begin{tabular}{lcccc}
\hline
Dataset & Tiny & Small & Normal & Large \\
\hline
Anti-UAV-RGBT IR (T\textsubscript{1} src) & 0.1\% & 3.3\% & 28.2\% & 68.4\% \\
Anti-UAV410 train (T\textsubscript{1} tgt) & 2.5\% & 27.5\% & 26.8\% & 43.2\% \\
CST train (T\textsubscript{2})             & 20.5\% & 73.6\% & 5.9\% & $<$0.01\% \\
\hline
\end{tabular}
\end{table}

The shift from T\textsubscript{1} (43.2\% large) to T\textsubscript{2}
(essentially 0\% large) represents a near-complete inversion of the
scale distribution. Furthermore, the elevated occlusion rate in CST
(14--18\% vs.\ 1--2\% in T\textsubscript{1} datasets) indicates a
substantially higher frequency of Thermal Crossover events, directly
relevant to RQ3 and the Evidential Uncertainty module.

% ------------------------------------------------------------
\subsection*{7. Data Processing Pipeline}

Following the dataset audit, a full data processing pipeline was
designed and submitted to the Snellius SLURM scheduler as a
dependency-chained job sequence. The pipeline covers both datasets
that require frame extraction and the annotation format conversions
required before training.

\subsubsection*{7.1 Rationale for Frame Extraction}

The initial strategy was to pre-extract all video frames to individual
JPEG files on disk before training. This decision was motivated by
four considerations:

Random access during training. Deep learning training requires
randomly sampling frames across all sequences in each mini-batch.
Retrieving a specific frame from an \texttt{.mp4} file requires
seeking through the video from the nearest keyframe, which introduces
unpredictable latency. Pre-extracted frames allow instant random
access via a single file read, regardless of temporal position in the
sequence.

Format consistency across datasets. Anti-UAV410 and CST
Anti-UAV are already distributed as pre-extracted JPEG frames. A
unified frame-on-disk representation would allow a single PyTorch
\texttt{Dataset} class to serve all four datasets without
branching on input format.

Decoupling data preparation from training. Pre-extracting
frames separates the data preparation phase from the training phase.
This is particularly important on shared HPC infrastructure where
GPU compute nodes should not be burdened with video decoding overhead
during time-limited GPU jobs.

Reproducibility. Extracted frames are immutable snapshots,
eliminating any dependency on codec versions or OpenCV build flags
that could otherwise affect how frames are decoded at different
points in time.

\subsubsection*{7.2 SLURM Job Chain}

Four jobs were submitted as a dependency chain, with Jobs 1a and 1b
running in parallel:

\begin{enumerate}
    \item Job 1a --- \texttt{extract\_antiuav\_rgbt}: Extracts
          frames from all infrared and visible \texttt{.mp4} video files
          across the train, val, and test splits of Anti-UAV-RGBT
          (318 sequences, both modalities). Includes a temporal alignment
          verification step confirming IR and RGB frame counts match per
          sequence.
    \item Job 1b --- \texttt{extract\_ard100}: Extracts frames
          from all ARD100 train and test \texttt{.mp4} videos (100
          sequences). Frame naming is matched exactly to the XML
          annotation filenames to ensure correct pairing.
    \item Job 2 --- \texttt{prepare\_ard100} (starts after Job
          1b): Creates the official train/test split
          \texttt{ImageSets/Main/} files using the predefined YOLOMG
          sequence split (65 train, 35 test), converts all Pascal VOC
          XML annotations to YOLO \texttt{.txt} format, and generates
          image path lists for the training framework.
    \item Job 3 --- \texttt{gen\_masks} (starts after Job 2):
          Generates pixel-level motion difference maps (\texttt{mask32})
          for all 100 ARD100 sequences using the \texttt{FD5\_mask}
          algorithm. These motion maps constitute YOLOMG's second input
          stream alongside the RGB frames.
\end{enumerate}

\subsubsection*{7.3 Anti-UAV JSON to YOLO Conversion (Attempted)}

A conversion script (\texttt{json2yolo.py}) was written to convert the
JSON annotations of Anti-UAV-RGBT and Anti-UAV410 into YOLO
\texttt{.txt} format, normalising bounding boxes from absolute pixel
coordinates $[x, y, w, h]$ to normalised $[x_c, y_c, w, h]$ using the
confirmed sensor resolution of $640 \times 512$ pixels.

% ------------------------------------------------------------
\subsection*{8. Infrastructure Roadblock: Inode Quota}

Both the frame extraction pipeline and the JSON-to-YOLO conversion
were terminated by the same underlying constraint: a per-project inode
quota on \texttt{/projects/prjs2041/}. An inode is a filesystem entry
representing a single file or directory; every extracted frame and
every per-frame \texttt{.txt} label file consumes one inode.

\subsubsection*{8.1 Scale of the Problem}

The proposed approach would have generated approximately:

\begin{itemize}
    \item Frame files: Anti-UAV-RGBT ($318$ sequences $\times$
          ${\sim}1{,}000$ frames $\times$ 2 modalities) $+$ ARD100
          ($100$ sequences $\times$ ${\sim}2{,}000$ frames)
          $\approx 836{,}000$ JPEG files.
    \item Label files: One \texttt{.txt} per frame across
          Anti-UAV-RGBT and Anti-UAV410
          $\approx 800{,}000$ additional files.
    \item Total: ${\sim}1.6$ million new filesystem inodes,
          which exceeds the project directory quota.
\end{itemize}

The extraction job crashed at sequence 27/160 of the Anti-UAV-RGBT
train split (\texttt{OSError: [Errno 122] Disk quota exceeded}), and
the label conversion crashed at sequence 128/160.

\subsubsection*{8.2 Revised Strategy: Native Format Reading via PyTorch Dataset Classes}

The inode constraint motivates a fundamental redesign of the data
loading strategy. Rather than converting all annotation formats to a
common on-disk representation, the revised approach reads each
dataset in its native format directly inside a PyTorch
\texttt{Dataset} class. No intermediate files are written to disk.
Table~\ref{tab:dataset_loading} summarises how each dataset is handled.

\begin{table}[h]
\centering
\caption{Native-format loading strategy per dataset.}
\label{tab:dataset_loading}
\resizebox{\columnwidth}{!}{%
\begin{tabular}{llll}
\hline
Dataset & Frame source & Annotation source & Extra files written \\
\hline
Anti-UAV-RGBT  & \texttt{cv2.VideoCapture} on \texttt{.mp4} & \texttt{infrared.json} / \texttt{visible.json} & None \\
Anti-UAV410    & \texttt{cv2.imread} on pre-extracted JPGs  & \texttt{IR\_label.json}                        & None \\
ARD100         & \texttt{cv2.VideoCapture} on \texttt{.mp4} & Pascal VOC XML per frame                       & None \\
CST Anti-UAV   & \texttt{cv2.imread} on pre-extracted JPGs  & \texttt{gt.txt} + \texttt{IR\_label.json}      & None \\
\hline
\end{tabular}%
}
\end{table}

Why this works.
Each \texttt{Dataset} class encapsulates all format-specific logic
internally: video seeking, JSON parsing, XML parsing, bounding box
normalisation, and the \texttt{exist} flag check. The rest of the
training pipeline — the DataLoader, the model, the loss function —
sees a uniform interface of \texttt{(image\_tensor, label\_tensor)}
pairs regardless of which dataset is being read. This is preferable
to the conversion approach even beyond the inode constraint, because
it removes a brittle preprocessing stage and keeps the original
annotation files as the single source of truth.

What this achieves.
The four Dataset classes will serve as the foundation for the entire
experimental pipeline. The Teacher-Student UDA framework (T\textsubscript{1})
requires a paired loader that returns synchronised RGB and IR frames
from Anti-UAV-RGBT — this is only cleanly expressible as a Dataset
class, not as a static file conversion. The Continual Learning
curriculum ($T_1 \to T_2$) requires switching between datasets
mid-training, which is straightforward when each dataset is a
self-contained class. The Forgetting Measure evaluation requires
replaying samples from earlier tasks, which maps naturally onto
the rehearsal buffer reading from Dataset instances.

Performance.
The loading overhead of on-the-fly video decoding is mitigated by:
(i) using multiple PyTorch DataLoader worker processes to prefetch
batches in parallel; (ii) caching the current sequence's
\texttt{VideoCapture} handle across consecutive frames; and (iii)
running data loading asynchronously on CPU while the GPU executes the
forward and backward passes.

% ------------------------------------------------------------
\subsection*{9. Next Steps}

\begin{enumerate}
    \item PyTorch Dataset classes: Implement a unified dataset
          interface for all four annotation formats using on-the-fly
          video reading for Anti-UAV-RGBT and ARD100, and direct image
          reading for Anti-UAV410 and CST. This constitutes the primary
          engineering task for the next session.
    \item Motion mask generation for ARD100: Adapt
          \texttt{generate\_mask5\_fixed.py} to generate motion
          difference maps directly into memory during training, or
          pre-compute and store them as a single compressed archive
          per sequence to minimise inode usage.
    \item Baseline training run: Establish the source-domain
          RGB baseline (mAP on Anti-UAV-RGBT) as the starting point
          for the Forgetting Measure calculation.
\end{enumerate}
